\section{Introduction}
\label{sec:intro}

Fork-join queueing systems model the parallel execution of tasks that must
all complete before a job can proceed.
A job arrives, \emph{forks} into $n$ tasks dispatched to $n$ parallel servers,
and \emph{joins}---departs---only when the last task finishes.
This synchronization constraint makes fork-join queues substantially harder
to analyze than independent queues, because the job response time is the
\emph{maximum} of the task sojourn times, and those times are stochastically
dependent.

We study the canonical case of $n = 2$ servers with Poisson arrivals at
rate $\lambda$ and exponential service times at rates $\mu_1$ and $\mu_2$,
with $\mu_1 \geq \mu_2$ without loss of generality.
The primary performance measure is the mean job response time:
\begin{equation}
  T = \E[\max(R_1, R_2)],
\end{equation}
where $R_i$ is the sojourn time (waiting plus service) of a task at server~$i$.
The system is stable when $\lambda < \mumin \equiv \mu_2$.

For the \emph{homogeneous} case $\mu_1 = \mu_2 = \mu$, Nelson and
Tantawi~\citep{Nelson1988} derived the exact closed-form result
\begin{equation}
  \label{eq:nelson-tantawi}
  T_2^{\mathrm{hom}} = \frac{12 - \rho}{8(\mu - \lambda)},
  \qquad \rho = \frac{\lambda}{\mu},
\end{equation}
using a light-heavy traffic interpolation argument.
For \emph{heterogeneous} servers ($\mu_1 \neq \mu_2$), the exact joint
distribution of $(R_1, R_2)$ is not available in closed form.
Flatto and Hahn~\citep{Flatto1984,Flatto1985} obtained the exact generating
function via an elliptic function parametrization, but extracting the mean
response time requires solving a boundary value problem and does not yield
a tractable expression.
Bounds, heavy-traffic asymptotics, and distributional
approximations have been developed~\citep{Baccelli1989,Kemper2012,Rizk2015,Mohanty2024},
but no simple closed-form approximation covering the full range of loads and
heterogeneity levels has previously been established.

In this paper we propose and analyze two complementary closed-form
approximations:
\begin{enumerate}
\item \textbf{$\Tul$ (upper-lower bound interpolation):}
      A convex combination of the independent upper bound $\Tub$ and the
      bottleneck lower bound $\Tbot$, weighted by the average utilization.
      It is guaranteed to lie within the bounds and is exact for the
      homogeneous case.
\item \textbf{$\Tlh$ (light-heavy traffic interpolation):}
      A rational approximation in the bottleneck utilization, derived from the
      Reiman--Simon framework~\citep{Reiman1988} as described
      in~\citep{Squillante2026}.
      It matches the exact light-traffic value and a heavy-traffic factor
      $h(r) = 1 + \tfrac{3}{8}\,r^{-\beta}$ that smoothly interpolates
      between the homogeneous limit ($h = 11/8$) and the bottleneck limit
      ($h \to 1$) as a function of the heterogeneity ratio $r = \mu_1/\mu_2$.
\end{enumerate}

Both approximations reduce exactly to the Nelson--Tantawi result
(\ref{eq:nelson-tantawi}) in the homogeneous case, exhibit the correct
singularity structure in heavy traffic, and agree with simulation to within
$2.5\%$ across all tested scenarios.

The remainder of the paper is organized as follows.
Section~\ref{sec:bounds} reviews known upper and lower bounds.
Sections~\ref{sec:approx-ul} and~\ref{sec:approx-lh} derive the two
approximations.
Section~\ref{sec:validation} presents simulation validation.
Section~\ref{sec:discussion} discusses the relationship between the two
methods, their limitations, and connections to prior work.
Section~\ref{sec:conclusion} concludes.
